\documentclass[aspectratio=169]{beamer}
\usetheme{Madrid}
\usecolortheme{default}
\usepackage[utf8]{inputenc}
\usepackage[spanish,es-nodecimaldot]{babel}
\usepackage{amsmath}
\usepackage{amssymb}
\usepackage{booktabs}
\usepackage{multirow}
\usepackage{graphicx}

\definecolor{unsblue}{RGB}{30,60,120}
\setbeamercolor{structure}{fg=unsblue}
\setbeamercolor{palette primary}{bg=unsblue,fg=white}
\setbeamercolor{palette secondary}{bg=unsblue!80,fg=white}
\setbeamercolor{palette tertiary}{bg=unsblue!60,fg=white}

\setbeamertemplate{navigation symbols}{}
\setbeamertemplate{footline}[frame number]

\title[MAD1 - Clase 28]{Métodos de Análisis de Datos 1}
\subtitle{Clase 28: Comunicación de Resultados y Ética}
\author{Prof. José}
\institute{Departamento de Matemática \\ Universidad Nacional del Sur}
\date{}

\begin{document}

\begin{frame}
\titlepage
\end{frame}

\begin{frame}{Contenidos de la Clase}
\tableofcontents
\end{frame}

\section{Introducción}

\begin{frame}{¿Por qué Comunicación y Ética?}
\textbf{El mejor análisis es inútil si:}
\begin{itemize}
\item No podemos explicarlo claramente
\item La audiencia no lo entiende
\item No genera confianza
\item Causa daño no intencional
\end{itemize}

\vspace{0.3cm}

\begin{block}{Realidad profesional}
Como científicos de datos, pasamos tanto tiempo comunicando y justificando nuestro trabajo como haciendo análisis técnico.
\end{block}

\vspace{0.3cm}

\textbf{Esta clase cubre:}
\begin{itemize}
\item Principios de comunicación efectiva
\item Reportes reproducibles (Jupyter, R Markdown)
\item Ética: sesgos, fairness, privacidad, responsabilidad
\end{itemize}
\end{frame}

\section{Comunicación de Resultados}

\begin{frame}{Principio 1: Conocer a la Audiencia}
\begin{columns}
\begin{column}{0.5\textwidth}
\textbf{Audiencia Técnica}\\
(científicos de datos, ingenieros)

\begin{itemize}
\item Detalles técnicos bienvenidos
\item Hiperparámetros, métricas específicas
\item Gráficos diagnósticos
\item Metodología rigurosa
\item Reproducibilidad
\end{itemize}
\end{column}

\begin{column}{0.5\textwidth}
\textbf{Audiencia de Negocio}\\
(gerentes, ejecutivos)

\begin{itemize}
\item Enfoque en impacto y ROI
\item Decisiones accionables
\item Lenguaje simple
\item Visualizaciones intuitivas
\item Resultados clave
\end{itemize}
\end{column}
\end{columns}

\vspace{0.4cm}

\begin{alertblock}{Error común}
Presentar detalles técnicos a audiencia no técnica (o viceversa)
\end{alertblock}
\end{frame}

\begin{frame}{Principio 2: Estructura Clara}
\textbf{Estructura de reporte típica:}

\begin{enumerate}
\item \textbf{Resumen ejecutivo} (1 página): problema, metodología, resultados, recomendaciones
\item \textbf{Introducción:} contexto, objetivos, preguntas
\item \textbf{Datos:} fuentes, tamaño, calidad, limitaciones
\item \textbf{Metodología:} exploración, features, modelos, validación
\item \textbf{Resultados:} métricas, comparaciones, hallazgos
\item \textbf{Conclusiones:} interpretación, implicaciones
\item \textbf{Limitaciones:} supuestos, sesgos, incertidumbre
\item \textbf{Próximos pasos:} mejoras, investigación futura
\item \textbf{Apéndices:} detalles técnicos, código, tablas
\end{enumerate}
\end{frame}

\begin{frame}{Principio 3: Visualizaciones Efectivas}
\textbf{Reglas de oro:}

\begin{enumerate}
\item \textbf{Un gráfico, un mensaje}
\begin{itemize}
\item No sobrecargar con información
\end{itemize}

\item \textbf{Claridad total}
\begin{itemize}
\item Ejes etiquetados, leyendas, títulos descriptivos
\end{itemize}

\item \textbf{Honestidad}
\begin{itemize}
\item No distorsionar (escalas apropiadas, barras desde cero)
\end{itemize}

\item \textbf{Accesibilidad}
\begin{itemize}
\item Colores para daltónicos, fuente legible
\end{itemize}

\item \textbf{Consistencia}
\begin{itemize}
\item Mismo estilo en todo el reporte
\end{itemize}
\end{enumerate}
\end{frame}

\begin{frame}{Elegir el Gráfico Apropiado}
\begin{table}
\small
\begin{tabular}{ll}
\toprule
\textbf{Objetivo} & \textbf{Gráfico} \\
\midrule
Comparar categorías & Barras, columnas \\
Evolución temporal & Líneas \\
Distribución & Histogramas, boxplots, violines \\
Relación entre dos variables & Scatter plots \\
Proporciones & Pie charts (con moderación), treemaps \\
Múltiples variables & Heatmaps, parallel coordinates \\
Composición & Áreas apiladas, treemaps \\
Ranking & Barras horizontales ordenadas \\
\bottomrule
\end{tabular}
\end{table}

\vspace{0.3cm}

\begin{alertblock}{Evitar}
Gráficos 3D innecesarios, efectos visuales que distorsionan, demasiados colores
\end{alertblock}
\end{frame}

\begin{frame}{Principio 4: Narrar una Historia}
\textbf{Los datos por sí solos no son convincentes}

\vspace{0.3cm}

\textbf{Estructura narrativa:}
\begin{enumerate}
\item \textbf{Contexto:} ¿por qué es importante?
\item \textbf{Problema:} ¿qué pregunta respondemos?
\item \textbf{Viaje:} ¿qué descubrimos en el camino?
\item \textbf{Clímax:} hallazgo principal
\item \textbf{Resolución:} ¿qué hacemos con esta información?
\end{enumerate}

\vspace{0.3cm}

\textbf{Técnicas narrativas:}
\begin{itemize}
\item Comenzar con ejemplo concreto o anécdota
\item Usar analogías para conceptos complejos
\item Anticipar y responder objeciones
\item Terminar con llamado a la acción claro
\end{itemize}
\end{frame}

\begin{frame}{Reportar Incertidumbre}
\textbf{Todo modelo tiene incertidumbre}

\vspace{0.3cm}

\textbf{Comunicarla honestamente:}
\begin{itemize}
\item Intervalos de confianza, no solo puntos
\item Errores de predicción contextualizados (RMSE, MAE)
\item Limitaciones de los datos (sesgos, faltantes)
\item Supuestos del modelo y su razonabilidad
\end{itemize}

\vspace{0.3cm}

\begin{exampleblock}{Ejemplo de buena práctica}
``El modelo predice una tasa de churn del 15\% con IC 95\%: [12\%, 18\%]. Esta estimación asume comportamiento futuro similar al histórico. Cambios en el mercado o productos podrían afectarla.''
\end{exampleblock}
\end{frame}

\begin{frame}{Lenguaje Claro y Preciso}
\begin{columns}
\begin{column}{0.5\textwidth}
\textbf{Evitar:}
\begin{itemize}
\item Jerga innecesaria
\item Siglas no definidas
\item Lenguaje vago
\item ``El modelo funciona bien''
\item Causalidad cuando solo hay correlación
\end{itemize}
\end{column}

\begin{column}{0.5\textwidth}
\textbf{Preferir:}
\begin{itemize}
\item Términos simples
\item Definir en primera mención
\item Cuantificación precisa
\item ``85\% de accuracy''
\item ``X está asociado con Y''
\end{itemize}
\end{column}
\end{columns}

\vspace{0.4cm}

\begin{alertblock}{Regla de oro}
Si tu abuela no lo entiende, necesita simplificarse (excepto para audiencia técnica especializada).
\end{alertblock}
\end{frame}

\section{Reportes Reproducibles}

\begin{frame}{¿Qué es Reproducibilidad?}
\begin{block}{Definición}
Un análisis es reproducible si otra persona puede ejecutarlo y obtener los mismos resultados.
\end{block}

\vspace{0.3cm}

\textbf{¿Por qué es importante?}
\begin{itemize}
\item \textbf{Verificación:} detectar errores
\item \textbf{Confianza:} transparencia genera credibilidad
\item \textbf{Colaboración:} otros pueden extender tu trabajo
\item \textbf{Ciencia:} resultados deben ser verificables
\item \textbf{Auditoría:} cumplimiento regulatorio
\end{itemize}

\vspace{0.3cm}

\textbf{Componentes de reproducibilidad:}
\begin{itemize}
\item Código versionado (Git)
\item Datos documentados
\item Dependencias especificadas (requirements.txt, renv)
\item Instrucciones claras
\end{itemize}
\end{frame}

\begin{frame}{Jupyter Notebooks}
\textbf{Combina código, visualizaciones y narrativa}

\vspace{0.3cm}

\textbf{Ventajas:}
\begin{itemize}
\item Interactividad: ejecutar celda por celda
\item Integración: código + texto + ecuaciones + imágenes
\item Compartible: formato JSON, versionable con Git
\item Exportable: HTML, PDF, presentaciones
\item Popular en industria y academia
\end{itemize}

\vspace{0.3cm}

\textbf{Desventajas:}
\begin{itemize}
\item Puede desordenarse fácilmente
\item Outputs grandes dificultan versionado
\item Tentación de código no ordenado
\end{itemize}
\end{frame}

\begin{frame}{Jupyter: Buenas Prácticas}
\begin{enumerate}
\item \textbf{Orden lógico:} celdas ejecutables secuencialmente de arriba a abajo

\item \textbf{Restart \& Run All:} verificar que funciona desde cero

\item \textbf{Markdown generoso:} explicar qué hace cada sección

\item \textbf{Imports al inicio:} todas las librerías arriba

\item \textbf{Configuración temprana:} seeds, parámetros, paths

\item \textbf{Control de versiones:} commits frecuentes

\item \textbf{Limpiar outputs:} antes de commitear (o usar nbstripout)

\item \textbf{Gestión de dependencias:} requirements.txt o environment.yml
\end{enumerate}
\end{frame}

\begin{frame}{Jupyter: Estructura Recomendada}
\begin{enumerate}
\item \textbf{Título y resumen}
\item \textbf{Imports:} todas las librerías
\item \textbf{Configuración:} seeds, parámetros
\item \textbf{Carga de datos:} con verificaciones
\item \textbf{Exploración:} EDA
\item \textbf{Preprocesamiento:} limpieza, transformaciones
\item \textbf{Modelado:} entrenamiento, validación
\item \textbf{Evaluación:} métricas, gráficos
\item \textbf{Conclusiones:} hallazgos principales
\item \textbf{Próximos pasos:} mejoras futuras
\end{enumerate}

\vspace{0.3cm}

\textbf{Exportar:}
\begin{itemize}
\item HTML: \texttt{jupyter nbconvert --to html notebook.ipynb}
\item PDF: \texttt{jupyter nbconvert --to pdf notebook.ipynb}
\end{itemize}
\end{frame}

\begin{frame}[fragile]{R Markdown}
\textbf{Combina código R, narrativa y análisis}

\vspace{0.3cm}

\textbf{Estructura de archivo .Rmd:}

\begin{enumerate}
\item \textbf{Encabezado YAML:} metadatos
\begin{verbatim}
---
title: "Análisis Q1 2026"
author: "José Bavio"
output: html_document
---
\end{verbatim}

\item \textbf{Texto Markdown:} narrativa

\item \textbf{Chunks de código R:} entre \texttt{```\{r\}} y \texttt{```}
\end{enumerate}
\end{frame}

\begin{frame}[fragile]{R Markdown: Opciones de Chunks}
Controlan cómo se ejecuta y muestra el código:

\vspace{0.2cm}

\begin{itemize}
\item \texttt{echo=FALSE}: no mostrar código, solo output
\item \texttt{eval=FALSE}: mostrar código pero no ejecutar
\item \texttt{include=FALSE}: ejecutar pero no mostrar nada
\item \texttt{message=FALSE, warning=FALSE}: suprimir mensajes
\item \texttt{fig.width=8, fig.height=6}: tamaño de gráficos
\item \texttt{cache=TRUE}: cachear resultados (útil si cálculos lentos)
\end{itemize}

\vspace{0.3cm}

\textbf{Ejemplo:}
\begin{verbatim}
```{r analisis-descriptivo, echo=FALSE, fig.width=10}
summary(datos)
ggplot(datos, aes(x=variable)) + geom_histogram()
```
\end{verbatim}
\end{frame}

\begin{frame}[fragile]{R Markdown: Formatos de Salida}
\textbf{HTML (interactivo):}
\begin{verbatim}
output: 
  html_document:
    toc: true
    toc_float: true
    theme: cosmo
\end{verbatim}

\textbf{PDF (vía LaTeX):}
\begin{verbatim}
output: pdf_document
\end{verbatim}

\textbf{Word:}
\begin{verbatim}
output: word_document
\end{verbatim}

\textbf{Presentaciones:}
\begin{verbatim}
output: ioslides_presentation
\end{verbatim}
\end{frame}

\begin{frame}[fragile]{R Markdown: Reportes Parametrizados}
\textbf{Crear reportes automatizados con parámetros variables}

\vspace{0.2cm}

\textbf{Definir parámetros:}
\begin{verbatim}
---
title: "Reporte de Ventas"
params:
  region: "Norte"
  año: 2026
---
\end{verbatim}

\textbf{Usar en el reporte:}
\begin{verbatim}
## Ventas en `r params$region` - `r params$año`

```{r}
datos_filtrados <- datos %>%
  filter(region == params$region, año == params$año)
```
\end{verbatim}

\textbf{Renderizar con diferentes parámetros:}
\begin{verbatim}
rmarkdown::render("reporte.Rmd", 
                  params = list(region="Sur", año=2025))
\end{verbatim}
\end{frame}

\begin{frame}{Control de Versiones con Git}
\textbf{Esencial para reproducibilidad}

\vspace{0.3cm}

\textbf{Beneficios:}
\begin{itemize}
\item Historial completo de cambios
\item Colaboración simultánea
\item Branching: experimentar sin romper nada
\item Rollback: volver a versiones anteriores
\item Transparencia y auditoría
\end{itemize}

\vspace{0.3cm}

\textbf{Workflow básico:}
\begin{enumerate}
\item \texttt{git init} (inicializar repo)
\item \texttt{git add archivo.ipynb} (agregar cambios)
\item \texttt{git commit -m "Mensaje descriptivo"}
\item \texttt{git push origin main} (subir a remoto)
\end{enumerate}
\end{frame}

\begin{frame}[fragile]{Gestión de Dependencias}
\textbf{Documentar qué librerías y versiones se usaron}

\vspace{0.3cm}

\textbf{Python - requirements.txt:}
\begin{verbatim}
pandas==2.0.1
numpy==1.24.3
scikit-learn==1.2.2
matplotlib==3.7.1
\end{verbatim}

Generar: \texttt{pip freeze > requirements.txt}

Instalar: \texttt{pip install -r requirements.txt}

\vspace{0.3cm}

\textbf{R - renv:}
\begin{verbatim}
renv::init()      # Inicializar
renv::snapshot()  # Guardar estado
renv::restore()   # Reinstalar desde snapshot
\end{verbatim}
\end{frame}

\section{Ética en Análisis de Datos}

\begin{frame}{¿Por qué Ética?}
\begin{block}{Poder conlleva responsabilidad}
Los modelos de ML tienen poder para afectar vidas: contratación, créditos, justicia penal, salud.
\end{block}

\vspace{0.3cm}

\textbf{Consecuencias de análisis no ético:}
\begin{itemize}
\item Discriminación sistemática contra grupos vulnerables
\item Violación de privacidad
\item Pérdida de confianza en sistemas algorítmicos
\item Daño a reputación profesional y organizacional
\item Consecuencias legales
\end{itemize}

\vspace{0.3cm}

\textbf{Como científicos de datos, somos responsables} de considerar el impacto ético de nuestro trabajo.
\end{frame}

\begin{frame}{Sesgo en Datos y Modelos}
\textbf{Fuentes de sesgo:}

\begin{enumerate}
\item \textbf{Sesgo de selección:} datos no representativos
\begin{itemize}
\item Ejemplo: entrenar modelo de contratación con datos históricos sesgados
\end{itemize}

\item \textbf{Sesgo de medición:} errores sistemáticos en recolección
\begin{itemize}
\item Ejemplo: encuestas telefónicas excluyen a quienes no tienen teléfono
\end{itemize}

\item \textbf{Sesgo algorítmico:} el algoritmo amplifica sesgos
\begin{itemize}
\item Ejemplo: sistema de recomendación refuerza burbujas de filtro
\end{itemize}

\item \textbf{Sesgo de confirmación:} analistas buscan datos que confirmen creencias
\begin{itemize}
\item Ejemplo: cherry-picking resultados favorables
\end{itemize}
\end{enumerate}
\end{frame}

\begin{frame}{Caso de Estudio: COMPAS}
\textbf{Correctional Offender Management Profiling for Alternative Sanctions}

\vspace{0.3cm}

\textbf{Sistema:} Evalúa riesgo de reincidencia en sistema judicial de EEUU

\vspace{0.3cm}

\textbf{Problema detectado:}
\begin{itemize}
\item Predecía falsamente alta reincidencia para acusados afroamericanos
\item Falsos positivos: 45\% para afroamericanos vs 23\% para blancos
\item Impacto: decisiones de libertad condicional y sentencias afectadas
\end{itemize}

\vspace{0.3cm}

\textbf{Lecciones:}
\begin{itemize}
\item Modelos pueden perpetuar y amplificar discriminación histórica
\item Auditorías externas son necesarias
\item Alto impacto requiere altos estándares de fairness
\item Transparencia es crucial
\end{itemize}
\end{frame}

\begin{frame}{Mitigación de Sesgo}
\textbf{Estrategias:}

\begin{enumerate}
\item \textbf{Auditoría de datos:}
\begin{itemize}
\item ¿Son representativos? ¿Hay subrepresentación?
\item ¿Cómo se recolectaron?
\end{itemize}

\item \textbf{Análisis de equidad:}
\begin{itemize}
\item Evaluar performance por subgrupo (género, raza, edad)
\item Métricas de fairness: disparate impact, equalized odds
\end{itemize}

\item \textbf{Técnicas de mitigación:}
\begin{itemize}
\item Pre-procesamiento: rebalancear datos
\item In-processing: regularización de fairness
\item Post-procesamiento: ajustar umbrales por grupo
\end{itemize}

\item \textbf{Diversidad en equipos:}
\begin{itemize}
\item Perspectivas diversas identifican sesgos
\item Incluir stakeholders afectados
\end{itemize}
\end{enumerate}
\end{frame}

\begin{frame}{Fairness: Múltiples Definiciones}
\textbf{No hay una definición única de fairness}

\vspace{0.3cm}

\textbf{1. Demographic Parity:}
\begin{itemize}
\item Misma proporción de resultados positivos para todos los grupos
\item $P(\hat{Y}=1 | A=a) = P(\hat{Y}=1 | A=b)$
\end{itemize}

\vspace{0.2cm}

\textbf{2. Equalized Odds:}
\begin{itemize}
\item Mismas tasas de VP y FP entre grupos
\item $P(\hat{Y}=1 | Y=1, A=a) = P(\hat{Y}=1 | Y=1, A=b)$
\end{itemize}

\vspace{0.2cm}

\textbf{3. Predictive Parity:}
\begin{itemize}
\item Mismo valor predictivo positivo
\item $P(Y=1 | \hat{Y}=1, A=a) = P(Y=1 | \hat{Y}=1, A=b)$
\end{itemize}

\vspace{0.3cm}

\begin{alertblock}{Imposibilidad}
En general, no se pueden satisfacer múltiples definiciones simultáneamente. Debemos elegir según el contexto.
\end{alertblock}
\end{frame}

\begin{frame}{Privacidad y Protección de Datos}
\textbf{Principios fundamentales:}

\begin{enumerate}
\item \textbf{Minimización:} recolectar solo lo necesario
\item \textbf{Propósito específico:} usar datos solo para lo autorizado
\item \textbf{Consentimiento informado:} las personas deben saber qué y cómo
\item \textbf{Transparencia:} derecho a saber qué datos se tienen
\item \textbf{Derecho al olvido:} poder eliminar datos personales
\item \textbf{Seguridad:} proteger contra acceso no autorizado
\end{enumerate}

\vspace{0.3cm}

\textbf{Regulaciones:}
\begin{itemize}
\item \textbf{GDPR (UE):} protección de datos, derecho a explicación
\item \textbf{CCPA (California):} derechos de consumidores
\item \textbf{Ley 25.326 (Argentina):} protección de datos personales
\end{itemize}
\end{frame}

\begin{frame}{Técnicas de Privacidad}
\textbf{1. Anonimización:}
\begin{itemize}
\item Eliminar identificadores directos (nombre, DNI)
\item Problema: reidentificación mediante quasi-identificadores
\end{itemize}

\vspace{0.2cm}

\textbf{2. K-anonymity:}
\begin{itemize}
\item Cada registro indistinguible de al menos $k-1$ otros
\item Generalización y supresión de atributos
\end{itemize}

\vspace{0.2cm}

\textbf{3. Differential Privacy:}
\begin{itemize}
\item Agregar ruido aleatorio a resultados
\item Garantía matemática de privacidad
\item Trade-off privacidad vs utilidad
\end{itemize}

\vspace{0.2cm}

\textbf{4. Federated Learning:}
\begin{itemize}
\item Entrenar sin centralizar datos
\item Datos permanecen en dispositivos locales
\end{itemize}
\end{frame}

\begin{frame}{Responsabilidad en Uso de Modelos}
\textbf{Decisiones Humanas vs Algorítmicas}

\vspace{0.3cm}

\textbf{Cuándo usar algoritmos:}
\begin{itemize}
\item Decisiones repetitivas, alto volumen
\item Criterios objetivos claramente definidos
\item Datos de alta calidad
\item Consecuencias moderadas de errores
\end{itemize}

\vspace{0.3cm}

\textbf{Cuándo mantener humanos:}
\begin{itemize}
\item Alto impacto individual (salud, justicia)
\item Contextos complejos que requieren juicio
\item Situaciones novedosas
\item Casos donde explicación y empatía importan
\end{itemize}

\vspace{0.3cm}

\textbf{Enfoque recomendado:} Human-in-the-loop (modelo recomienda, humano decide)
\end{frame}

\begin{frame}{Auditoría y Rendición de Cuentas}
\textbf{Prácticas esenciales:}

\begin{itemize}
\item \textbf{Auditorías regulares:} revisar performance, sesgo, errores

\item \textbf{Trazabilidad:} registrar decisiones y justificaciones

\item \textbf{Responsables identificados:} quién responde por el sistema

\item \textbf{Mecanismos de apelación:} personas pueden disputar decisiones

\item \textbf{Evaluaciones de impacto:} antes del despliegue, evaluar riesgos

\item \textbf{Documentación completa:} Model Cards, Datasheets for Datasets

\item \textbf{Monitoreo continuo:} detectar degradación o sesgos emergentes
\end{itemize}
\end{frame}

\begin{frame}{Preguntas Éticas a Hacerse}
\textbf{Antes de empezar un proyecto:}

\begin{enumerate}
\item ¿Tenemos derecho a estos datos? ¿Hubo consentimiento?
\item ¿Para qué se usarán los resultados? ¿Podría usarse para daño?
\item ¿Quién se beneficia y quién podría ser perjudicado?
\item ¿Los datos son representativos? ¿Hay sesgos conocidos?
\item ¿Qué tan confiables son las predicciones? ¿Margen de error?
\item ¿Hay alternativas no algorítmicas? ¿Son preferibles?
\item ¿Quién será responsable si algo sale mal?
\item ¿Podemos explicar las decisiones? ¿A quién, cómo?
\end{enumerate}

\vspace{0.3cm}

\begin{alertblock}{Señal de alerta}
Si no puedes responder estas preguntas satisfactoriamente, reconsidera el proyecto.
\end{alertblock}
\end{frame}

\begin{frame}{Red Flags Éticas}
\textbf{Señales de alerta que indican problemas:}

\begin{itemize}
\item Cliente solicita ``eliminar variables sensibles pero mantener performance''
\item Objetivos vagos o cambiantes
\item Presión para reportar solo resultados favorables
\item Falta de transparencia sobre uso final
\item Datos obtenidos de manera cuestionable
\item Stakeholders afectados no consultados
\item Consecuencias asimétricas (FP mucho peores que FN, o viceversa)
\item Resistencia a auditorías externas
\item Falta de mecanismos de apelación
\end{itemize}

\vspace{0.3cm}

\textbf{Si ves estas señales:} cuestiona, documenta, escala, o considera rechazar el proyecto.
\end{frame}

\begin{frame}{Código de Ética Profesional}
\textbf{Principios a seguir como científicos de datos:}

\begin{enumerate}
\item \textbf{Integridad:} reportar hallazgos honestamente

\item \textbf{Transparencia:} documentar limitaciones, supuestos

\item \textbf{Responsabilidad:} considerar impacto de nuestro trabajo

\item \textbf{Equidad:} esforzarse por sistemas justos

\item \textbf{Privacidad:} proteger datos personales

\item \textbf{Competencia:} reconocer límites de nuestra expertise

\item \textbf{Beneficencia:} buscar que nuestro trabajo beneficie a la sociedad
\end{enumerate}

\vspace{0.3cm}

\textbf{Referencia:} ACM Code of Ethics (Association for Computing Machinery)
\end{frame}

\section{Resumen}

\begin{frame}{Resumen de la Clase}
\textbf{Comunicación:}
\begin{itemize}
\item Conocer a la audiencia, estructura clara
\item Visualizaciones efectivas y honestas
\item Narrar una historia, reportar incertidumbre
\end{itemize}

\vspace{0.3cm}

\textbf{Reportes reproducibles:}
\begin{itemize}
\item Jupyter notebooks: interactivo, integrado
\item R Markdown: automatizable, múltiples formatos
\item Control de versiones (Git), gestión de dependencias
\end{itemize}

\vspace{0.3cm}

\textbf{Ética:}
\begin{itemize}
\item Reconocer y mitigar sesgos
\item Múltiples definiciones de fairness
\item Proteger privacidad (anonimización, differential privacy)
\item Responsabilidad: auditorías, transparencia, rendición de cuentas
\item Preguntas críticas antes de cada proyecto
\end{itemize}
\end{frame}

\begin{frame}{Cierre del Curso}
\textbf{Recorrido completo del ciclo de análisis de datos:}

\begin{enumerate}
\item \textbf{U0:} Introducción, entorno, reproducibilidad
\item \textbf{U1:} Obtención y preparación de datos
\item \textbf{U2:} Exploración y visualización
\item \textbf{U3:} Modelado estadístico (regresión, regularización)
\item \textbf{U4:} Machine learning supervisado (clasificación)
\item \textbf{U5:} Producción, interpretabilidad, comunicación, ética
\end{enumerate}

\vspace{0.3cm}

\textbf{Próximos pasos:}
\begin{itemize}
\item Proyecto integrador final
\item Aplicar todo lo aprendido
\item Construir portfolio profesional
\item Seguir aprendiendo (¡esto nunca termina!)
\end{itemize}
\end{frame}

\begin{frame}{Mensaje Final}
\begin{center}
\Large
\textbf{La ciencia de datos es poder.}

\vspace{0.5cm}

\textbf{Usemos ese poder sabiamente.}

\vspace{1cm}

\normalsize
Como científicos de datos, tenemos la responsabilidad de:
\begin{itemize}
\item Crear modelos que funcionen en el mundo real
\item Explicar nuestros hallazgos clara y honestamente
\item Documentar para reproducibilidad y auditoría
\item Considerar el impacto en las personas
\item Abogar por el uso ético de datos y algoritmos
\end{itemize}

\vspace{0.5cm}

\textbf{¡Éxitos en su carrera profesional!}
\end{center}
\end{frame}

\end{document}
